
\subsubsection*{4.1.1}
\phantomsection
\label{sga3:VIIB:4.1.1}
\footnote{\label{sga3:VIIB:nde:134}The editors corrected the
original, which gave the inclusion
\[
  (\widehat X/\widehat S)^{(p)}
  \subset
  \widehat{X^{(p)}}/\widehat S
\]
instead of the reverse inclusion. They also note that this
paragraph is not used below.}

Let \(S\) be a scheme over the prime field \(\mathbb F_p\), and
let \(\eta:X\to S\) be an \(S\)-scheme. Write
\(\operatorname{fr}(S)\) for the ``absolute'' Frobenius morphism
of \(S\): it induces the identity on the underlying topological
space and sends every section \(f\) of the structure sheaf to
\(f^p\); cf.\ VII\textsubscript{A}, \SGARefLink{sga3:VIIA:4.1}{4.1}. Let
\[
  X^{(p)}=X\times_{\operatorname{fr}(S)}S
\]
be the corresponding \(S\)-scheme
\textup{(VII\textsubscript{A}, loc.\ cit.)}; thus its structural
morphism to \(S\) is \(\operatorname{fr}(S)\circ\eta\).

Let \(\widehat S\) be the formal scheme whose underlying
topological space is the set of closed points of \(S\), endowed
with the discrete topology, and whose local ring at a closed point
\(s\) is the separated completion
\[
  \mathcal O_{\widehat S,s}
  =
  \widehat{\mathcal O}_{S,s}
\]
of \(\mathcal O_{S,s}\) for the linear topology defined by ideals
of finite colength \textup{(cf.\ \SGARefLink{sga3:VIIB:1.2.6}{1.2.6})}. Its affine algebra is
therefore
\[
  k=\mathcal A(\widehat S)
  =
  \prod_{s\in S_{\mathrm{cl}}}
  \widehat{\mathcal O}_{S,s}.
\]

Recall \textup{(loc.\ cit.)} that
\(\widehat X/\widehat S\) denotes the formal \(k\)-variety
consisting of those points \(x\in X\) for which
\[
  [\kappa(x):\kappa(s)]<\infty,
\]
where \(s\) is the image of \(x\) in \(S\). Its local ring at
\(x\) is the separated completion
\[
  \mathcal O_{\widehat X/\widehat S,x}
  =
  \widehat{\mathcal O}_{X,x}
\]
of \(\mathcal O_{X,x}\) for the linear topology defined by those
ideals \(I\) such that \(\mathcal O_{X,x}/I\) has finite length as
an \(\mathcal O_{X,x}\)-module, and hence also as an
\(\mathcal O_{S,s}\)-module. By base change we can therefore form
the formal \(k\)-variety
\[
  (\widehat X/\widehat S)^{(p)}
  =
  (\widehat X/\widehat S)
  \widehat\otimes_{\pi}k.
\]

We can also consider the formal \(k\)-variety
\(\widehat{X^{(p)}}/\widehat S\). Its underlying set consists of
the points \(x\in X^{(p)}=X\) such that, if \(s\) denotes the image
of \(x\) in \(S\), the composite
\begin{center}
  \phantomsection
  \label{sga3:VIIB:4.1.1:residue-field-frobenius-diagram}
  \(
\kappa(s)\xrightarrow{\ \operatorname{fr}\ }\kappa(s)
\xrightarrow{\ \eta^\natural\ }\kappa(x)
\)

\end{center}
makes \(\kappa(x)\) a finite extension of \(\kappa(s)\). In that
case the same is true of
\[
  \eta^\natural:\kappa(s)\longrightarrow\kappa(x),
\]
so \(x\) is a point of \(\widehat X/\widehat S\), and
\[
  \mathcal O_{\widehat{X^{(p)}}/\widehat S,x}
  =
  \widehat{\mathcal O}_{X,x}
  \widehat\otimes_{\pi}
  \widehat{\mathcal O}_{S,s}
  =
  \widehat{\mathcal O}_{X^{(p)},x}.
\]
The second equality follows because the construction
\(X\mapsto\widehat X/\widehat S\) commutes with products
\textup{(cf.\ \SGARefLink{sga3:VIIB:1.2.6}{1.2.6})}.

It follows that
\[
  \widehat{X^{(p)}}/\widehat S
  \ \subset\
  (\widehat X/\widehat S)^{(p)}
\]
as a formal subvariety. Equality holds if and only if, for every
closed point \(s\) of \(S\), the field \(\kappa(s)\) has finite
degree over \(\kappa(s)^p\). This is the case, for example, if
\(S\) is locally of finite type over a field \(\kappa\) such that
\[
  [\kappa:\kappa^p]<\infty.
\]
